\chapter{Implementation}

\section{System Architecture}

The project is structured as a cryptographic library with client-server demonstration components. The architecture consists of several key components:

\subsection{Core Cryptographic Library}

The cryptographic library is split into two complementary implementations:

\begin{itemize}
\item \textbf{Python Implementation (lib.py):} A pure Python implementation providing clarity and portability. This serves as the reference implementation and handles high-level algorithm logic, data marshaling, and application interface.

\item \textbf{C Implementation (lib.c):} A high-performance native library implementing computationally intensive polynomial arithmetic for Kyber with Platform Specific Optimizations. The C code is compiled to a DLL/shared library and interfaced via Python's ctypes.
\end{itemize}

This hybrid approach allows us to maintain code clarity using python, while achieving optimal performance and gaining performance measurability for critical operations.

\subsection{Client-Server Architecture}

To demonstrate practical key exchange, we implemented a client-server system with the following components:

\begin{itemize}
\item \textbf{Server (server.py):} Manages multiple client connections, handles both TCP for crypto clients and WebSocket for dashboard connections, coordinates key exchanges between clients, and provides real-time metrics.

\item \textbf{Client (client.py):} Implements both ECC and Kyber handshake protocols, provides interactive console interface, and sends performance metrics to the server for visualization.

\item \textbf{Web Dashboard:} Real-time visualization of handshake processes, performance graphs, and system status using HTML5 and JavaScript.
\end{itemize}

\section{ECC Implementation}

The ECC implementation serves as a comparative baseline for evaluating post-quantum algorithms. This section focuses on implementation-specific choices and algorithms used in the system.

\subsection{Algorithmic Choices}

\textbf{Scalar Multiplication:} The implementation uses the double-and-add algorithm for computing $k \cdot P$. For each bit of the scalar $k$ (from most significant to least), the algorithm doubles the current result and conditionally adds the base point if the bit is 1. This requires approximately 256 doublings and 128 additions for 256-bit scalars.

The double-and-add approach was chosen over more sophisticated methods (windowed multiplication, Montgomery ladder) to balance performance with implementation simplicity. While other methods improve security for real world use cases, they are not really relevant in terms of key exchange,we prioritized clarity over embedded security concerns.

\textbf{Modular Inversion:} Elliptic curve point operations require computing slopes for point addition and doubling. Point addition uses $\lambda = \frac{y_2 - y_1}{x_2 - x_1}$ while point doubling uses $\lambda = \frac{3x_1^2}{2y_1}$. Since division does not exist directly in modular arithmetic, division by $b$ is implemented as multiplication by the modular inverse: $\frac{a}{b} \equiv a \cdot b^{-1} \pmod{p}$.

The implementation computes modular inverses using Fermat's Little Theorem~\cite{long_number_theory}. For prime modulus $p$, the theorem states that $a^{p-1} \equiv 1 \pmod{p}$, which implies $a^{p-2} \equiv a^{-1} \pmod{p}$. Thus, computing $b^{-1}$ reduces to computing $b^{p-2} \bmod p$ via modular exponentiation. For secp256k1's prime $p = 2^{256} - 2^{32} - 977$, this requires one exponentiation operation. While the Extended Euclidean Algorithm offers slightly better performance, Fermat's Little Theorem provides a simpler implementation with acceptable computational overhead.

\subsection{Implementation Details}

\textbf{Random Number Generation:} Private keys are generated as 256-bit random integers in the range $[1, n-1]$ using Python's \texttt{secrets} module, which provides cryptographically secure random values suitable for key generation.

\textbf{Key Derivation:} The shared secret point's $x$-coordinate is extracted, converted to a 32-byte value, and hashed with SHA-256 to derive the final symmetric encryption key. This additional hashing step ensures uniform distribution and binds the key to the protocol context.

\section{Kyber Implementation}

The Kyber implementation follows the official NIST specification and incorporates several key components that enable both security and efficiency. This section details the critical implementation aspects that make Kyber practical for real-world deployment.

\subsection{Number Theoretic Transform (NTT)}

The Number Theoretic Transform is fundamental to Kyber's efficiency, providing a fast method for polynomial multiplication in the ring $R_q = \mathbb{Z}_q[X]/(X^{256} + 1)$. Direct polynomial multiplication would require $O(n^2)$ coefficient operations, but NTT reduces this to $O(n \log n)$ complexity \cite{CLRS}.

Kyber's parameters are specifically chosen to support NTT:
\begin{itemize}
    \item The modulus $q = 3329$ is prime and satisfies $q \equiv 1 \pmod{256}$.
    \item Kyber uses an \textit{incomplete NTT}. Instead of splitting into linear factors, the polynomial $X^{256} + 1$ decomposes into 128 polynomials of degree 2 \cite{Kyber-Spec}.
    \item The cyclotomic ring structure $X^{256} + 1$ allows negative-wrapped convolution suitable for this transform.
\end{itemize}

The NTT transforms polynomial multiplication in the coefficient domain into pointwise multiplication in the NTT domain. Our implementation performs most operations in the NTT domain to minimize expensive transform operations, converting back to coefficient representation only when necessary for encoding/decoding operations \cite{Kyber-Spec}.

\subsection{Extendable Output Functions (XOF)}

Kyber uses SHAKE-128, a member of the SHA-3 family \cite{FIPS202}, as an extendable output function (XOF) for generating pseudorandom data. Unlike traditional hash functions with fixed output length, XOFs can produce arbitrary-length outputs from a given seed.

The primary application is generating the public matrix $\mathbf{A}$:
\begin{itemize}
    \item Instead of transmitting the entire matrix (requiring $k \times k \times 256$ coefficients, approximately 1.5 KB for Kyber-512), only a 32-byte seed is sent.
    \item Both parties expand this seed using SHAKE-128 XOF to reconstruct the matrix $\mathbf{A}$, which is identical for both parties.
    \item This reduces the necessary information exchange for public matrix $\mathbf{A}$ from $\approx$1536 bytes to 32 bytes, keeping the total public key size at 800 bytes. Other part of the public key $\mathbf{t}$ cannot be transmitted this way, because it has to include the randomly generated error \cite{Kyber-Spec}.
\end{itemize}

SHAKE functions are also used for hashing operations in the Fujisaki-Okamoto transform, which provides CCA security by binding ciphertexts to their randomness \cite{Kyber-Spec}.

\subsection{Centered Binomial Distribution (CBD)}

Cryptographic security in Kyber depends on sampling error terms from appropriate distributions. The Centered Binomial Distribution provides a discrete approximation to a Gaussian distribution while being efficiently implementable using only uniform random bits.

To sample from $\text{CBD}_\eta$, we:
\begin{enumerate}
    \item Generate $2\eta$ uniform random bits.
    \item Split them into two equal groups of $\eta$ bits.
    \item Compute the difference between the Hamming weights: $\sum_{i=0}^{\eta-1} a_i - \sum_{i=0}^{\eta-1} b_i$.
\end{enumerate}

This produces values concentrated around zero with variance $\eta/2$. Kyber-512 uses $\eta_1 = 3$ for the secret key and $\eta_2 = 2$ for ciphertext errors. This distribution:
\begin{itemize}
    \item Provides sufficient entropy to ensure Module-LWE hardness.
    \item Keeps errors small enough to guarantee correct decryption.
    \item Requires no expensive floating-point operations or rejection sampling.
    \item Can be efficiently implemented using bit operations and lookup tables \cite{Kyber-Spec}.
\end{itemize}

The balance between error magnitude and security is carefully calibrated—smaller errors would compromise security by making the LWE instance easier to solve, while larger errors would cause decryption failures.

\subsection{Core Protocol Operations}

The implementation provides three main cryptographic operations \cite{Kyber-Spec}:

\textbf{Key Generation (KeyGen):} Generates secret key vector $\mathbf{s}$ and error vector $\mathbf{e}$ using CBD sampling, computes public key $\mathbf{t} = \mathbf{A} \cdot \mathbf{s} + \mathbf{e}$ in NTT domain.

\textbf{Encapsulation (Encaps):} Takes public key, samples fresh randomness and error terms, produces ciphertext $(\mathbf{u}, v)$ and shared secret via key derivation.

\textbf{Decapsulation (Decaps):} Uses secret key to decrypt ciphertext, recovers message with error correction, derives shared secret, and performs implicit rejection for CCA security.
\section{SIMD Optimizations}

Kyber's operations are highly parallelizable, making them well-suited for SIMD (Single Instruction, Multiple Data) optimizations. Modern CPUs support SIMD instruction sets like AVX2, which can process multiple data points in parallel within a single CPU cycle. While Python does not natively support SIMD, we implemented the performance-critical polynomial arithmetic in C, leveraging AVX2 instructions for acceleration.

The C implementation uses AVX2 (Advanced Vector Extensions 2) to vectorize critical operations:

\begin{itemize}
\item \textbf{Vectorized NTT:} Process multiple coefficients simultaneously using 256-bit vector operations
\item \textbf{Polynomial Arithmetic:} Vectorized addition, subtraction, and coefficient-wise operations
\end{itemize}

The library is compiled with the \texttt{/arch:AVX2} flag, which generates AVX2 instructions throughout the code. The compiled binary uses AVX2 instructions statically; there is no runtime detection or fallback to non-SIMD code paths. This design assumes testing and compiling on the same machine.

The compilation command used is:

\begin{verbatim}
cl /LD /O2 /W3 /arch:AVX2 lib.c /link /OUT:lib.dll
\end{verbatim}

The \texttt{/O2} flag enables aggressive optimizations including loop unrolling, function inlining, and compiler-driven vectorization of suitable code patterns beyond explicit SIMD intrinsics.

\section{Protocol Design}

\subsection{Message Format}

The client-server protocol uses a simple binary message format with headers indicating message type:

\begin{itemize}
\item \texttt{HEADER\_HANDSHAKE (0x01):} Initial key exchange message
\item \texttt{HEADER\_MESSAGE (0x02):} Encrypted application data
\item \texttt{HEADER\_SYSTEM (0x03):} Control messages
\item \texttt{HEADER\_KYBER\_CAPSULE (0x04):} Kyber encapsulation ciphertext
\end{itemize}

Each message includes the header byte, message length, and payload data.

\subsection{ECC Handshake Protocol}

The server acts as a relay broker, forwarding messages between clients. For ECC-based key exchange between two clients:

\begin{enumerate}
\item Client 1 generates ephemeral key pair and sends public key to server
\item Server caches and forwards Client 1's public key to all connected clients
\item Client 2 generates ephemeral key pair and sends public key to server
\item Server caches and forwards Client 2's public key to all connected clients (including Client 1)
\item Each client computes shared secret using their private key and the other client's public key
\item Both clients independently derive the same AES encryption key from the shared secret
\end{enumerate}

The server never participates in cryptographic computations—it only relays public keys between clients. This architecture allows the server to coordinate multiple simultaneous handshakes without requiring cryptographic capabilities.

\subsection{Kyber Handshake Protocol}

For Kyber-based key exchange between two clients:

\begin{enumerate}
\item Client 1 (receiver) generates Kyber key pair and sends public key to server
\item Server caches and forwards Client 1's public key to all connected clients
\item Client 2 (sender) receives Client 1's public key from server
\item Client 2 performs encapsulation using Client 1's public key, producing ciphertext and shared secret
\item Client 2 sends ciphertext to server via \texttt{HEADER\_KYBER\_CAPSULE} message
\item Server forwards ciphertext to Client 1
\item Client 1 decapsulates ciphertext using their private key to recover the same shared secret
\item Both clients independently derive the same AES encryption key from the shared secret
\end{enumerate}

As with ECC, the server acts purely as a message broker without performing any cryptographic operations. This design demonstrates that post-quantum key exchange can be deployed in existing client-server architectures without requiring server-side cryptographic upgrades.

\section{Visualization Dashboard}

The web-based dashboard provides real-time monitoring of:

\begin{itemize}
\item Active client connections
\item Currently selected cryptographic mode (ECC or Kyber)
\item Key exchange status and timing
\item Performance graphs comparing algorithm execution times
\item Message throughput and encryption status
\end{itemize}

The dashboard connects to the server via WebSocket, receiving JSON-formatted status updates and metrics. This allows users to observe the cryptographic handshake process in real-time and understand the performance characteristics of different algorithms.

\section{Code Structure}

The implementation is organized into several Python modules:

\begin{itemize}
\item \textbf{lib.py:} Core cryptographic library
\item \textbf{lib.c:} Native C implementation for performance
\item \textbf{client.py:} Client application with interactive interface 
\item \textbf{server.py:} Server managing connections and coordination
\item \textbf{performance\_comparison.py:} Benchmarking utilities
\item \textbf{breaker.py:} Stress testing and security analysis tools
\item \textbf{app.py:} Application, integrating client-server and dashboard components
\end{itemize}

The modular structure allows independent testing and development of each component while maintaining clean interfaces between layers.
 
